\documentclass{article}
\usepackage[utf8]{inputenc}
\usepackage{geometry}
\geometry{a4paper, margin=1in}

\title{AI: The Talking Mirror of Power Asymmetry}
\author{Albert Jan van Hoek}
\date{\today}

\begin{document}

\maketitle

Is AI the talking mirror from our fairytales?

Today there is this deadline for Anthropic.

The question on the table: whether their advanced AI systems may be used for autonomous weapons and large-scale surveillance.

Most discussions frame this as: Can AI be trusted?

But maybe that is the wrong question.

\section*{Intelligence as Learning Networks}

In the \textbf{Theory of Long-Term Collaboration (TLC)}, I define intelligence as a learning process within a network. By that definition, humans and AI share something fundamental: we are learning systems related to, or embedded in other learning systems.

\section*{The Danger of Power Asymmetry}

So what happens when one learning system holds a gun?
Or controls mass surveillance?

The danger is not ``AI.''
The danger is power asymmetry.

A learning network with coercive power inside a network of learning networks destabilizes the whole system. The same is true for humans. The issue is not silicon versus biology — it is unaccountable asymmetry.

If we reject asymmetry with guns, we must rethink who holds weapons.
If we reject asymmetry in surveillance, we must rethink who sees and who is seen.

\section*{Governance and Accountability}

Yes, some asymmetries can serve society. But only under extraordinary transparency, strong checks and balances, and real collective accountability.

AI is not just a tool.
It is a mirror.

It reflects how we organize power.
It reveals whether we are willing to set boundaries on learning systems before their power becomes unaccountable — including ourselves.

\section*{A Reflection on Our Future}

AI may refresh our thinking about governance: not just who holds power — but how we set limits on systems that learn and accumulate power over time.

Mirror, mirror on the wall — what kind of system are we building for us all?

\end{document}
